\BeleskeID{03-s018}
% element: 03-s018:e01
% element: 03-s018:e02
% element: 03-s018:e03
% element: 03-s018:d01

Iz koderskog dela sledi
\[
 A_1=I_3+\overline{I_3}I_2=I_3+I_2.
\]
Za donji bit važi
\[
\begin{aligned}
 A_0&=I_3+\overline{I_3}\overline{I_2}I_1
      =I_3+\overline{I_2}I_1,\\
 GS&=I_3+I_2+I_1+I_0.
\end{aligned}
\]
U prva dva izraza upotrebljen je identitet \(X+\overline X Y=X+Y\). GS ostaje jedan čim postoji bilo koji aktivan ulaz, bez obzira na to koji od njih ima prednost.

Inverzija \(I_2\) na putanji od \(I_1\) čuva prioritet: aktivni \(I_2\) blokira doprinos \(I_1\) donjem bitu. Ako je aktivan \(I_3\), njegov direktan doprinos postavlja oba adresna bita na jedan. Ulaz \(I_0\) menja samo GS, jer njegov indeks 0 ima oba adresna bita nula. Kompaktnija šema daje istu funkcionalnu tabelu, ali ne izdvaja sve međusignale prioriteta kao prethodna realizacija.
